\chapter{Testing and Evaluation}

\section{Testing Strategy}
Testing a quantum virtual machine is fundamentally different from testing a traditional web service. In classical software, arithmetic produces deterministic integers or discrete strings. In quantum simulation, arithmetic operates on complex amplitude vectors in Hilbert spaces where numerical rounding errors accumulate across sequential matrix products. A test cannot simply assert \texttt{state == expected}; it must evaluate unitary invariants, check probability normalization ($\sum |\alpha_i|^2 = 1.0$), and allow for floating-point tolerance ($\epsilon \sim 10^{-7}$).

We organized our testing strategy into four practical layers:
\begin{itemize}[leftmargin=*]
    \item \textbf{Unit Testing:} Validating individual compiler passes and math kernels (e.g., verifying that a Hadamard gate on $|0\rangle$ produces $[1/\sqrt{2}, 1/\sqrt{2}]^T$ within $10^{-7}$ tolerance).
    \item \textbf{Integration Testing:} Verifying data flow between adjacent pipeline modules (e.g., ensuring the Lark parser output correctly feeds the intermediate representation, which in turn feeds the transpiler).
    \item \textbf{Algorithm Correctness Testing:} Verifying end-to-end execution of canonical quantum algorithms (Bell state preparation, GHZ entanglement, Bernstein-Vazirani parity finding, and Grover's search) against known analytical baselines.
    \item \textbf{Cross-Framework Audit Benchmarking:} Comparing QVM's simulation outputs against industry-standard frameworks (Qiskit and Cirq) across a 20-algorithm corpus to verify semantic fidelity.
\end{itemize}

\textbf{Testing Framework and Organization:}
All tests are implemented using \texttt{pytest}. The suite spans 24 modules in the \texttt{tests/} directory, totaling 215 automated test cases:
\begin{itemize}[leftmargin=*]
    \item \textbf{Grammar and Parsers:} \texttt{test\_parser.py}, \texttt{test\_qasm\_parser.py}, \texttt{test\_qasm3\_extended.py}, \texttt{test\_qasm3\_loops.py}, \texttt{test\_qasm3\_shadow.py}, \texttt{test\_cirq\_parser.py}
    \item \textbf{Intermediate Representation:} \texttt{test\_ir.py}, \texttt{test\_json\_serialization.py}, \texttt{test\_qasm\_roundtrip.py}
    \item \textbf{Synthesis and Routing:} \texttt{test\_decomposer.py}, \texttt{test\_mc\_synthesis.py}, \texttt{test\_transpiler.py}, \texttt{test\_mps\_routing.py}
    \item \textbf{Simulation and Stress Suites:} \texttt{test\_simulator.py}, \texttt{test\_v03.py}, \texttt{test\_stress.py} (26 deep rotation and ansatz stress tests)
    \item \textbf{Framework Interoperability:} \texttt{test\_qiskit\_integration.py}, \texttt{test\_cirq\_integration.py}, \texttt{test\_interop\_roundtrip.py}, \texttt{test\_backend\_cross\_conversion.py}
    \item \textbf{Application and Visual:} \texttt{test\_api.py}, \texttt{test\_visual.py}, \texttt{test\_tutorials.py}
\end{itemize}

\begin{verbatim}
# Execute the automated test suite (215 tests)
python -m pytest tests/

# Execute cross-framework audit benchmarks
python -m benchmarks.run_audit
\end{verbatim}

\section{Unit and Integration Testing}
The automated test suite runs 215 automated tests covering both valid operational paths and deliberately malformed edge cases.

\subsection{Parser Tests}
The parser tests verify that input circuits are properly tokenized, validated, and translated into AST nodes.

\begin{table}[h!]
\centering
\caption{Parser Unit Tests}
\label{tab:parser_tests}
\begin{tabular}{|p{2cm}|p{5cm}|p{2cm}|p{3.5cm}|}
\hline
\textbf{Test ID} & \textbf{Test Description} & \textbf{Status} & \textbf{Expected Result} \\
\hline
UT-P-01 & Parse valid circuit with H, CX, RZ gates & Pass & \texttt{QuantumCircuit} with 3 operations \\
\hline
UT-P-02 & Parse empty circuit description & Pass & \texttt{QuantumCircuit} with 0 operations \\
\hline
UT-P-03 & Parse circuit missing 'name' field & Pass & \texttt{ValueError} raised \\
\hline
UT-P-04 & Parse circuit missing 'qubits' field & Pass & \texttt{ValueError} raised \\
\hline
UT-P-05 & Parse circuit with invalid qubit index & Pass & \texttt{ValueError} raised \\
\hline
\end{tabular}
\end{table}

\textbf{Coverage Analysis:} Parser tests cover 95\% of statements across \texttt{parser.py} and \texttt{qasm3\_parser.py}. Edge cases explicitly tested include undeclared classical registers, mismatched register dimensions, floating-point angle parameters, and nested conditional blocks.

\subsection{Simulator Tests}
Simulator tests check the mathematical accuracy of quantum state transitions.

\begin{table}[h!]
\centering
\caption{Simulator Unit Tests}
\label{tab:sim_tests}
\begin{tabular}{|p{2cm}|p{5cm}|p{2cm}|p{3.5cm}|}
\hline
\textbf{Test ID} & \textbf{Test Description} & \textbf{Status} & \textbf{Expected Result} \\
\hline
UT-S-01 & Apply Hadamard gate to $|0\rangle$ & Pass & $\frac{1}{\sqrt{2}}(|0\rangle + |1\rangle)$ \\
\hline
UT-S-02 & Apply Pauli-X gate to $|0\rangle$ & Pass & $|1\rangle$ \\
\hline
UT-S-03 & Apply RY($\pi$/2) gate to $|0\rangle$ & Pass & $\frac{1}{\sqrt{2}}(|0\rangle + |1\rangle)$ \\
\hline
UT-S-04 & Create Bell state (H + CX) & Pass & $\frac{1}{\sqrt{2}}(|00\rangle + |11\rangle)$ \\
\hline
UT-S-05 & Create GHZ state (H + 2 CX) & Pass & $\frac{1}{\sqrt{2}}(|000\rangle + |111\rangle)$ \\
\hline
UT-S-06 & Handle unsupported gate & Pass & \texttt{ValueError} raised \\
\hline
UT-S-07 & Handle invalid qubit count for gate & Pass & \texttt{ValueError} raised \\
\hline
UT-S-08 & Sample Bell state (2000 shots) & Pass & $\approx$50\% |00⟩, $\approx$50\% |11⟩ \\
\hline
UT-S-09 & Measurement collapse to basis state & Pass & Single basis state, norm=1 \\
\hline
\end{tabular}
\end{table}

\textbf{Validation Method:} Numerical assertions use \texttt{numpy.allclose()} with an absolute tolerance of $10^{-7}$. Stochastic measurement sampling runs across 2,000 shots, asserting that observed outcome frequencies fall within a $\pm 5\%$ confidence interval around the theoretical Born probabilities.

\subsection{Transpiler Tests}
Transpiler tests verify that routing algorithms insert valid SWAP operations without altering logical quantum semantics.

\begin{table}[h!]
\centering
\caption{Transpiler Unit Tests}
\label{tab:trans_tests}
\begin{tabular}{|p{2cm}|p{5cm}|p{2cm}|p{3.5cm}|}
\hline
\textbf{Test ID} & \textbf{Test Description} & \textbf{Status} & \textbf{Expected Result} \\
\hline
UT-T-01 & Transpile on fully connected architecture & Pass & No SWAP gates inserted \\
\hline
UT-T-02 & Transpile CX(0,2) on linear chain & Pass & SWAP gates inserted \\
\hline
UT-T-03 & Transpile on disconnected architecture & Pass & \texttt{RuntimeError} raised \\
\hline
UT-T-04 & Verify logical equivalence after transpilation & Pass & Statevectors match \\
\hline
UT-T-05 & SABRE reduces SWAPs vs Greedy & Pass & Fewer SWAP gates \\
\hline
\end{tabular}
\end{table}

\textbf{Semantic Invariance Verification:} The most important transpiler test asserts that compiling a circuit for a restricted coupling topology preserves its original quantum state. The test executes the unconstrained logical circuit, then transpiles and executes the physical circuit with inserted SWAPs, asserting that both final statevectors match within $10^{-7}$ tolerance.

\section{Integration Testing}
Integration tests evaluate interactions across multiple pipeline stages working together.

\begin{table}[h!]
\centering
\caption{Integration Tests}
\label{tab:integration_tests}
\begin{tabular}{|p{2cm}|p{6cm}|p{2cm}|p{2.5cm}|}
\hline
\textbf{Test ID} & \textbf{Test Description} & \textbf{Status} & \textbf{Components} \\
\hline
IT-01 & Parse QASM $\to$ Simulate $\to$ Visualize & Pass & Parser, Simulator, Visualizer \\
\hline
IT-02 & Parse JSON $\to$ Decompose $\to$ Simulate & Pass & Parser, Decomposer, Simulator \\
\hline
IT-03 & Parse $\to$ Transpile $\to$ Simulate & Pass & Parser, Transpiler, Simulator \\
\hline
IT-04 & CLI end-to-end execution & Pass & CLI, Parser, Simulator, Visualizer \\
\hline
IT-05 & API request $\to$ Execute $\to$ Response & Pass & API, Parser, Simulator \\
\hline
IT-06 & OpenQASM 3.0 with control flow & Pass & Parser, Simulator (labels/jumps) \\
\hline
\end{tabular}
\end{table}

\section{System Testing}
System tests validate full quantum algorithms against their mathematical specifications.

\begin{table}[h!]
\centering
\caption{Algorithm Verification Tests}
\label{tab:algo_tests}
\begin{tabular}{|p{2cm}|p{5cm}|p{2cm}|p{3.5cm}|}
\hline
\textbf{Test ID} & \textbf{Algorithm} & \textbf{Status} & \textbf{Verification Criteria} \\
\hline
ST-01 & Bell State Preparation & Pass & $|\psi\rangle = \frac{1}{\sqrt{2}}(|00\rangle + |11\rangle)$, 50/50 counts \\
\hline
ST-02 & GHZ-3 State Generation & Pass & $|\psi\rangle = \frac{1}{\sqrt{2}}(|000\rangle + |111\rangle)$ \\
\hline
ST-03 & Bernstein-Vazirani ($s=101$) & Pass & Measures hidden string $|101\rangle$ with $P=1.0$ \\
\hline
ST-04 & Grover Search (3 qubits, $w=101$) & Pass & Measured $|101\rangle$ with $P \ge 0.94$ \\
\hline
ST-05 & Quantum Teleportation & Pass & Teleports state $|\psi\rangle = \alpha|0\rangle + \beta|1\rangle$ \\
\hline
ST-06 & Superdense Coding & Pass & Transmits 2 classical bits via 1 qubit \\
\hline
ST-07 & VQE Ground State ($H_2$) & Pass & Converges within 1~mHa of exact FCI \\
\hline
ST-08 & QAOA MaxCut (4-node cycle) & Pass & Finds optimal partition cuts \\
\hline
\end{tabular}
\end{table}

\section{Performance and Benchmark Evaluation}

\subsection{Simulation Performance}
We benchmarked the dense statevector simulator across qubit counts from 2 to 14 qubits on an AMD Ryzen 7 processor with 16~GB RAM.

\begin{table}[h!]
\centering
\caption{Statevector Simulation Performance}
\label{tab:perf_statevec}
\begin{tabular}{|p{2.5cm}|p{3cm}|p{3.5cm}|p{3.5cm}|}
\hline
\textbf{Qubits ($n$)} & \textbf{Statevector Dimension} & \textbf{Memory Allocation} & \textbf{Runtime (50 Gates)} \\
\hline
2 & 4 & 64 B & 0.42 ms \\
\hline
4 & 16 & 256 B & 0.85 ms \\
\hline
6 & 64 & 1.0 KB & 1.94 ms \\
\hline
8 & 256 & 4.1 KB & 4.81 ms \\
\hline
10 & 1,024 & 16.4 KB & 18.2 ms \\
\hline
12 & 4,096 & 65.5 KB & 84.6 ms \\
\hline
14 & 16,384 & 262.1 KB & 382.1 ms \\
\hline
\end{tabular}
\end{table}

As shown in Table \ref{tab:perf_statevec}, memory grows strictly as $2^n \times 16$~bytes. Runtime remains below 100 milliseconds through 12 qubits, making interactive local experimentation snappy and responsive.

\subsection{Transpilation Performance}
Table \ref{tab:perf_transpiler} compares Greedy BFS against SABRE on a 5-qubit linear architecture ($Q_0 - Q_1 - Q_2 - Q_3 - Q_4$).

\begin{table}[h!]
\centering
\caption{Transpilation Performance Comparison}
\label{tab:perf_transpiler}
\begin{tabular}{|p{3.5cm}|p{2.5cm}|p{2.5cm}|p{3.5cm}|}
\hline
\textbf{Circuit Benchmark} & \textbf{Original CNOTs} & \textbf{Greedy SWAPs} & \textbf{SABRE SWAPs (Decay=0.5)} \\
\hline
Bell State & 1 & 0 & 0 \\
\hline
GHZ-3 State & 2 & 1 & 1 \\
\hline
Bernstein-Vazirani & 3 & 4 & 2 (50\% reduction) \\
\hline
Grover Search ($N=8$) & 8 & 12 & 7 (41.6\% reduction) \\
\hline
Random Circuit ($D=10$) & 15 & 22 & 14 (36.3\% reduction) \\
\hline
\end{tabular}
\end{table}

SABRE consistently achieves a 35\% to 50\% reduction in SWAP count for non-trivial circuits, confirming the value of lookahead evaluation over greedy choices.

\subsection{Matrix Product State (MPS) Scalability}
Table \ref{tab:perf_mps} documents MPS simulation throughput for 1D linear cluster states with bounded bond dimension $\chi = 16$.

\begin{table}[h!]
\centering
\caption{MPS Simulation Scalability}
\label{tab:perf_mps}
\begin{tabular}{|p{2.5cm}|p{3cm}|p{3.5cm}|p{3.5cm}|}
\hline
\textbf{Qubits ($n$)} & \textbf{Bond Dimension $\chi$} & \textbf{Truncation Error ($\Delta$)} & \textbf{Runtime (s)} \\
\hline
10 & 16 & $< 10^{-14}$ & 0.12 s \\
\hline
15 & 16 & $< 10^{-13}$ & 0.28 s \\
\hline
20 & 16 & $< 10^{-12}$ & 0.54 s \\
\hline
25 & 16 & $< 10^{-11}$ & 0.89 s \\
\hline
30 & 16 & $< 10^{-10}$ & 1.42 s \\
\hline
\end{tabular}
\end{table}

While dense statevector simulation exhausts 16~GB RAM around 30 qubits, the MPS simulator executes 30-qubit low-entanglement circuits in 1.42 seconds with negligible truncation error ($\Delta < 10^{-10}$).

\subsection{Cross-Framework Audit Corpus}
To rigorously evaluate QVM against industry standards, we ran an automated audit comparing QVM simulation outputs against Qiskit and Cirq across a corpus of 20 canonical circuits, summarized in Table \ref{tab:audit_corpus}.

\begin{table}[h!]
\centering
\caption{Cross-Framework Quantum Algorithm Audit Benchmark Corpus}
\label{tab:audit_corpus}
\begin{tabular}{|p{1.5cm}|p{4.5cm}|p{2cm}|p{4cm}|}
\hline
\textbf{ID} & \textbf{Algorithm / Benchmark} & \textbf{Qubits} & \textbf{Cross-Validation Status} \\
\hline
AUD-01 & Bell State $|\Phi^+\rangle$ & 2 & Match ($L_2 \text{ err} < 10^{-14}$) \\
\hline
AUD-02 & Bell State $|\Psi^-\rangle$ & 2 & Match ($L_2 \text{ err} < 10^{-14}$) \\
\hline
AUD-03 & GHZ 3-Qubit Entanglement & 3 & Match ($L_2 \text{ err} < 10^{-14}$) \\
\hline
AUD-04 & GHZ 5-Qubit Entanglement & 5 & Match ($L_2 \text{ err} < 10^{-14}$) \\
\hline
AUD-05 & Superdense Coding & 2 & Match ($L_2 \text{ err} < 10^{-14}$) \\
\hline
AUD-06 & Quantum Teleportation & 3 & Match ($L_2 \text{ err} < 10^{-14}$) \\
\hline
AUD-07 & Deutsch-Jozsa (Balanced) & 3 & Match ($L_2 \text{ err} < 10^{-14}$) \\
\hline
AUD-08 & Deutsch-Jozsa (Constant) & 3 & Match ($L_2 \text{ err} < 10^{-14}$) \\
\hline
AUD-09 & Bernstein-Vazirani ($s=101$) & 4 & Match ($L_2 \text{ err} < 10^{-14}$) \\
\hline
AUD-10 & Simon's Periodicity Algorithm & 4 & Match ($L_2 \text{ err} < 10^{-14}$) \\
\hline
AUD-11 & Grover Search ($N=4$, $w=11$) & 2 & Match ($L_2 \text{ err} < 10^{-14}$) \\
\hline
AUD-12 & Grover Search ($N=8$, $w=101$) & 3 & Match ($L_2 \text{ err} < 10^{-14}$) \\
\hline
AUD-13 & Quantum Fourier Transform (QFT) & 3 & Match ($L_2 \text{ err} < 10^{-13}$) \\
\hline
AUD-14 & Inverse QFT & 3 & Match ($L_2 \text{ err} < 10^{-13}$) \\
\hline
AUD-15 & Quantum Phase Estimation & 4 & Match ($L_2 \text{ err} < 10^{-13}$) \\
\hline
AUD-16 & Shor's Order Finding ($N=15$) & 5 & Match ($L_2 \text{ err} < 10^{-12}$) \\
\hline
AUD-17 & Random Clifford Benchmark & 4 & Match ($L_2 \text{ err} < 10^{-14}$) \\
\hline
AUD-18 & Multi-Controlled Toffoli (MCX) & 4 & Match ($L_2 \text{ err} < 10^{-14}$) \\
\hline
AUD-19 & VQE Molecular $H_2$ Ansatz & 2 & Match ($L_2 \text{ err} < 10^{-13}$) \\
\hline
AUD-20 & QAOA MaxCut (4-Node Cycle) & 4 & Match ($L_2 \text{ err} < 10^{-13}$) \\
\hline
\end{tabular}
\end{table}

All 20 algorithms produced identical statevectors and probability distributions between QVM, Qiskit, and Cirq, demonstrating semantic equivalence across frameworks.

\section{Security and Fault Tolerance Testing}
Because QVM accepts circuit code and JSON payloads from external users, we conducted defensive input validation tests.

\begin{table}[h!]
\centering
\caption{Security and Defensive Input Tests}
\label{tab:sec_tests}
\begin{tabular}{|p{2cm}|p{5cm}|p{2cm}|p{3.5cm}|}
\hline
\textbf{Test ID} & \textbf{Vulnerability Target} & \textbf{Status} & \textbf{Observed Defense Behavior} \\
\hline
SEC-01 & Directory Traversal (\texttt{../../etc/passwd}) & Pass & Rejected with \texttt{PathSecurityError} \\
\hline
SEC-02 & Infinite Loop via QASM3 \texttt{while} & Pass & Halted at \texttt{max\_ops = 100,000} \\
\hline
SEC-03 & Out-of-bounds Qubit Memory Exhaustion & Pass & Rejected before tensor allocation \\
\hline
SEC-04 & Malformed JSON Injection & Pass & Caught by FastAPI Pydantic validator \\
\hline
SEC-05 & Non-Unitary Matrix Application & Pass & Rejected during decomposition pass \\
\hline
\end{tabular}
\end{table}

\section{User Acceptance Testing (UAT)}
To evaluate system usability, five undergraduate computer science students and research assistants participated in user acceptance testing sessions.

\begin{table}[h!]
\centering
\caption{UAT Participants}
\label{tab:uat_participants}
\begin{tabular}{|p{2cm}|p{4.5cm}|p{6cm}|}
\hline
\textbf{ID} & \textbf{Role / Background} & \textbf{Prior Quantum Experience} \\
\hline
UAT-P1 & CS Final Year Student & Completed introductory Quantum Computing elective \\
\hline
UAT-P2 & CS Junior Student & Familiar with Python, beginner to quantum algorithms \\
\hline
UAT-P3 & Teaching Assistant & Experienced with Qiskit and scientific Python \\
\hline
UAT-P4 & Software Engineering Major & Web developer, evaluated REST API and dashboard \\
\hline
UAT-P5 & CS Honors Student & Researches compiler optimization techniques \\
\hline
\end{tabular}
\end{table}

Participants completed five standard workflow tasks: loading a circuit file via CLI, transpiling for a 5-qubit linear architecture, executing a circuit on the web dashboard, adding depolarizing noise, and reading the resulting probability histogram. All five participants completed the tasks successfully, awarding an average usability rating of 4.4 out of 5.0.

\section{Test Results Summary}

\begin{table}[h!]
\centering
\caption{Overall Test Results Summary}
\label{tab:overall_tests}
\begin{tabular}{|p{4.5cm}|p{2.5cm}|p{2.5cm}|p{2.5cm}|}
\hline
\textbf{Test Category} & \textbf{Total Test Cases} & \textbf{Passed} & \textbf{Success Rate} \\
\hline
Automated Pytest Unit \& IR Tests & 124 & 124 & 100\% \\
\hline
Synthesis \& Transpiler Suites & 48 & 48 & 100\% \\
\hline
Deep Rotation \& Stress Tests & 26 & 26 & 100\% \\
\hline
Framework Integration (Qiskit/Cirq) & 17 & 17 & 100\% \\
\hline
Cross-Framework Audit Corpus & 20 & 20 & 100\% \\
\hline
Security \& Input Validation & 5 & 5 & 100\% \\
\hline
User Acceptance Scenarios & 5 & 5 & 100\% \\
\hline
\textbf{Total Validations} & \textbf{245} & \textbf{245} & \textbf{100\%} \\
\hline
\end{tabular}
\end{table}

\textbf{Execution Record:}  
\begin{itemize}[leftmargin=*, noitemsep]
    \item \textbf{Automated Test Suite:} 215 tests passed, 1 skipped in 16.21 seconds.
    \item \textbf{Branch Coverage:} 95\% statement coverage across core compiler and simulator modules.
    \item \textbf{Numerical Tolerance:} Statevectors match analytical baselines within $\text{atol} = 10^{-7}$.
\end{itemize}

\section{Threats to Validity}

To maintain empirical and scientific rigor, we explicitly identify potential threats to the validity of our testing outcomes and benchmark conclusions:

\begin{itemize}[leftmargin=*]
    \item \textbf{Internal Validity (Floating-Point Precision and Drift):} Statevector and MPS simulations operate using IEEE 754 standard double-precision complex numbers (\texttt{numpy.complex128}). While single-qubit gate unitaries are theoretically orthonormal, repeated floating-point matrix multiplications across deep circuit chains ($>100$ gate operations) can introduce cumulative rounding errors. We mitigated this risk by enforcing periodic statevector norm renormalization and establishing strict validation thresholds ($\text{atol} = 10^{-7}$) during automated testing.
    \item \textbf{External Validity (Corpus Generalizability):} Our cross-framework audit evaluated 20 canonical quantum circuits (spanning superposition, entanglement, Grover search, and phase estimation). Although these circuits cover the foundational curriculum of quantum computing, highly irregular application-specific circuits (such as deep random circuit sampling or chemistry Hamiltonians with thousands of Pauli terms) may exhibit different SWAP insertion overheads and classical runtime profiles.
    \item \textbf{Construct Validity (Endianness and Bit-Ordering Alignment):} Different quantum computing frameworks adopt conflicting bit-ordering conventions. Qiskit indexes qubits in little-endian order ($q_0$ as the least significant bit), whereas OpenQASM 3.0 and QVM's internal statevector follow big-endian register indexing ($q_0$ as the most significant bit). Without rigorous conversion, automated comparisons can falsely report statevector mismatches. To guarantee construct validity, our test harness incorporated automated permutation matrices to invert register indexing before computing Euclidean state vector distances.
\end{itemize}

\section{Chapter Summary}
This chapter evaluated the correctness, numerical stability, performance, and security of the Quantum Virtual Machine. We detailed our testing methodology, explaining why quantum verification requires numerical tolerances rather than exact equality. We documented results across unit tests, pipeline integration tests, canonical quantum algorithms, and performance benchmarks. The statevector engine achieved sub-second execution up to 12 qubits, while the Matrix Product State engine demonstrated scalability to 30 qubits for low-entanglement states. SABRE lookahead routing demonstrated a 35\% to 50\% reduction in SWAP count over greedy routing. Finally, an audit across 20 algorithms confirmed numerical equivalence with Qiskit and Cirq. Chapter 7 concludes the report and discusses avenues for future research.
